% Part: incompleteness
% Chapter: arithmetization-syntax
% Section: introduction

\documentclass[../../../include/open-logic-section]{subfiles}

\begin{document}

\olfileid[hi]{inc}{art}{int}
\olsection{परिचय}

संगणनीयता और तर्कशास्त्र को जोड़ने के लिए हमें तार्किक वस्तुओं (चिह्नों,
पदों, !!{formula} और !!{derivation}), उन पर संक्रियाओं तथा उनके गुणों और
संबंधों की चर्चा ऐसे ढंग से करने का उपाय चाहिए जो संगणकीय उपचार के अनुकूल हो।
हम चिह्नों, चिह्न-अनुक्रमों और उनसे बनी अन्य वस्तुओं पर संगणनीय फलनों और
संबंधों पर सीधे विचार करके ऐसा कर सकते हैं। तार्किक वाक्यविन्यास की सभी
वस्तुएँ परिमित हैं और चिह्नों के !!a{enumerable} समुच्चयों से बनी हैं, इसलिए
संगणना के कुछ प्रतिरूपों में यह संभव है। किंतु संगणना के अन्य प्रतिरूप—जैसे
पुनरावर्ती फलन—संख्याओं तथा उनके संबंधों और फलनों तक सीमित हैं। इसके अतिरिक्त,
अंततः हम कुछ सिद्धांतों के भीतर, विशेषतः अंकगणित की भाषा में सूत्रित सिद्धांतों
में, वाक्यविन्यास को संभालना चाहते हैं। इन स्थितियों में वाक्यविन्यास का
\emph{अंकगणितीकरण} आवश्यक है; अर्थात् वाक्यविन्यासी वस्तुओं, उन पर संक्रियाओं
और उनके संबंधों को क्रमशः संख्याओं, अंकगणितीय फलनों और अंकगणितीय संबंधों के
रूप में निरूपित करना होगा। लाइबनित्स तक पीछे जाने वाला विचार वाक्यविन्यासी
वस्तुओं को संख्याएँ निर्दिष्ट करना है।

चिह्नों को उनके ``कूट'' के रूप में संख्याएँ देना अपेक्षाकृत सरल है। कुछ चिह्न
थोड़ी चुनौती देते हैं, क्योंकि उदाहरणतः अनंततः अनेक !!{variable} और प्रत्येक
स्थान-संख्या~$n$ के अनंततः अनेक !!{function} हैं। फिर भी चिह्नों को व्यवस्थित
ढंग से संख्याएँ देना संभव है ताकि, मान लें, $\Obj v_2$ और $\Obj v_3$ को भिन्न
कूट मिलें। चिह्न-अनुक्रम (जैसे पद और !!{formula}) बड़ी चुनौती हैं। किंतु यदि
हम संख्या-अनुक्रमों को शुद्ध अंकगणितीय ढंग से संभाल सकें (उदाहरणतः अनुक्रमों के
अभाज्य-घात कूटलेखन से), तो अलग-अलग चिह्नों के कूटलेखन को चिह्न-अनुक्रमों तक,
और फिर !!{formula} के अनुक्रमों या !!{derivation} जैसी अन्य व्यवस्थाओं तक
विस्तृत कर सकते हैं। यह विस्तृत कूटलेखन ``ग्योडेल (G\"odel) क्रमांकन''
कहलाता है। प्रत्येक पद, !!{formula} और !!{derivation} को एक ग्योडेल (G\"odel)
संख्या दी जाती है।

चिह्न-अनुक्रमों को उनके कूटों के अनुक्रमों के रूप में कूटित करके और अनुक्रमों
का ऐसा कूटलेखन-तंत्र चुनकर जिसे संगणनीय फलनों से संभाला जा सके, हम ग्योडेल (G\"odel)
 संख्याओं को भी संगणनीय फलनों से संभाल सकते हैं। व्यवहार में सभी
संबंधित फलन आदिम पुनरावर्ती होंगे। उदाहरणतः किसी अनुक्रम की लंबाई निकालना और
उसके कूट से उसका $i$-वाँ अवयव निकालना, दोनों आदिम पुनरावर्ती हैं। यदि अनुक्रम
को कूटित करने वाली संख्या, उदाहरणतः, !!a{formula}~$!A$ की ग्योडेल (G\"odel)
संख्या है, तो तुरंत दिखता है कि !!a{formula} की लंबाई और उस !!a{formula} में $i$-वें चिह्न
का कूट भी~$!A$ की ग्योडेल (G\"odel) संख्या से निकाला जा सकता है। यह सिद्ध
करना कुछ कठिन है कि, उदाहरणतः, किसी संख्या का सुगठित पद या सही
!!{derivation} की ग्योडेल (G\"odel) संख्या होना आदिम पुनरावर्ती गुण है। फिर भी
यह संभव है, क्योंकि संबंधित अनुक्रम—पद, !!{formula} और !!{derivation}—आगमनात्मक
रूप से परिभाषित हैं।

उदाहरण के लिए प्रतिस्थापन संक्रिया पर विचार करें। यदि $!A$ कोई सूत्र, $x$
कोई चर और $t$ कोई पद है, तो $\Subst{!A}{t}{x}$ वह परिणाम है जो~$!A$ में~$x$
की प्रत्येक मुक्त उपस्थिति को~$t$ से बदलने पर मिलता है। अब मान लें कि हमने
$!A$, $x$, $t$ को क्रमशः $k$, $l$ और $m$ ग्योडेल (G\"odel) संख्याएँ दी हैं।
यही योजना $\Subst{!A}{t}{x}$ को, मान लें,~$n$ ग्योडेल (G\"odel) संख्या देती
है। $k$, $l$ और $m$ को $n$ पर भेजने वाला यह प्रतिचित्रण प्रतिस्थापन संक्रिया
का अंकगणितीय अनुरूप है। जब प्रतिस्थापन संक्रिया $!A$, $x$, $t$ को
$\Subst{!A}{t}{x}$ पर भेजती है, तब अंकगणितीकृत प्रतिस्थापन फलन ग्योडेल (G\"odel)
 संख्याओं $k$, $l$, $m$ को ग्योडेल (G\"odel) संख्या~$n$ पर भेजता
है। हम देखेंगे कि यह फलन आदिम पुनरावर्ती है।

वाक्यविन्यास का अंकगणितीकरण केवल अमूर्त रुचि का विषय नहीं। मूलतः यह अपने आप
में एक अतुच्छ अंतर्दृष्टि थी कि अंकगणित की भाषा जैसी भाषाएँ, जिनमें भाषाओं के
विषय में ``बात करने'' की कोई व्यवस्था अंतर्निहित नहीं, फिर भी अभिव्यक्तियों
के जटिल गुणों को औपचारिक बना सकती हैं। इसके बाद यह पूछना केवल एक छोटा कदम है
कि इस भाषा का कोई सिद्धांत, जैसे पियानो अंकगणित, अपनी ही भाषा के विषय में
क्या \emph{सिद्ध} कर सकता है (उदाहरणतः क्या !!{sentence} सिद्धनीय या सत्य
हैं)। यह हमें ग्योडेल (G\"odel) के असिद्धनीयता संबंधी और टार्स्की के सत्य की
अपरिभाष्यता संबंधी प्रसिद्ध परिसीमन प्रमेयों तक ले जाता है। किंतु वाक्यविन्यास
के अंकगणितीकरण की युक्ति संगणनीयता सिद्धांत के कुछ महत्त्वपूर्ण परिणाम सिद्ध
करने के लिए भी आवश्यक है; उदाहरणतः सिद्धांतों की संगणकीय शक्ति या संगणनीयता
के विभिन्न प्रतिरूपों के संबंध के विषय में। वाक्यविन्यास का अंकगणितीकरण अन्य
वस्तुओं और गुणों के अंकगणितीकरण का प्रतिमान बनता है। उदाहरण के लिए, विन्यासों
और संगणनाओं (मान लें, ट्यूरिंग मशीनों की) का भी इसी प्रकार अंकगणितीकरण संभव
है। इससे एक प्रतिरूप (जैसे ट्यूरिंग मशीन) की संगणनाओं का दूसरे प्रतिरूप (जैसे
पुनरावर्ती फलनों) में अनुकरण संभव होता है।

\end{document}
